\section{Knowledge and Power}

\begin{quotex}
The fundamental being and expression [of Consciousness or the I] is Joy pulsating as Will-Power and manifesting Itself in an unspeakably sublime cosmic play. It is not a mere abstraction — a wilderness of Pure Being or Pure Nothing as some critics of Vedanta have imagined the abode of Reality to be. \flright{\textsc{John Woodroffe}, \emph{The World as Power} (p 320)}

\end{quotex}
Thomas Aquinas describes the effects of the Fall from the Primordial State as affecting faculties on both the spiritual and the psychic level, as represented in the following chart:

\begin{table}[h]
\centering\small
\begin{tabularx}{\textwidth}{llp{.3\textwidth}l}
\toprule
 &
\textbf{Effect} &
\textbf{Description} &
\textbf{Antidote}\\\midrule
Spirit &
Ignorance &
Intellect not ordered to the Logos &
Gnosis\\\cmidrule{2-4}
 &
Malice &
Will impotent to manifest the Logos &
Power\\\midrule
Soul &
Weakness &
Thumos not ordered to arduousness &
Trial by suffering\\\cmidrule{2-4}
 &
Concupiscence &
Eros not ordered to the delectable &
Trial by love\\\bottomrule
\end{tabularx}
\end{table}
Hence, the regeneration of the Primordial State is effected through Knowledge and Power.

\paragraph{Knowledge}
What do we know when we know it? From what has been described thus far:

\begin{itemize}
\item All our opinions, beliefs, theories are like so much straw 
\item What we have considered as obstacles are in reality privations, a demonstration of our own lack of power. 
\end{itemize}
Woodroffe describes the concept of privation as understood in Tantra:

\begin{quotex}
Consciousness [the I] \emph{seeming} to be unconsciousness, Joy \emph{seeming} to be indifference or pain, free Play \emph{seeming} to be necessity and determination. 

\end{quotex}
In this seemingness, the I gives up his power to that which ``seems to be" and doesn't recognize his own role in producing what he experiences:

\begin{quotex}
Consciousness, as the root Being-and-Becoming Power becomes real things … Thus, a block of stone is not ``matter" only: it is Consciousness as Joy and as Play — though the fact is veiled to ordinary minds. On the other hand, it is not an ``idea" or ``mental construct" only: it is Consciousness as Power constituting it as much and as active a reality as the experiencing and reacting mind is. 

\end{quotex}
The solution is, according to Woodroffe, increasing knowledge at ever deeper levels to the point of self-realization:

\begin{quotex}
There is need for the education and development of man's ``knowing instruments", giving him progressively higher and larger visions, through science and philosophy, through intuition and meditation, and finally revelation and realization. … \emph{Matter and Life are every whit as real and active as Mind}.

\end{quotex}
This is the crux of the issue. Verbosity, opinions, repetition do not and never will constitute knowledge. \textbf{To Know} — in the highest sense, viz., episteme, gnosis, wisdom — is \textbf{to be}. Hence, we do not know the Idea until it has become Real for us as Life and as Matter. This is the Tantric teaching of the I as Power.

\paragraph{Power}
In the \emph{Individual and the Becoming of the World}, Evola describes the deepening process of making ideas real:

\begin{enumerate}
\item First, there is the simple and pallid image of the idea 
\item Make it vivid in the imagination 
\item Perceive it exteriorly as a subjective hallucination 
\item Convince other centres of consciousness to perceive it as a collective hallucination 
\item Develop an ever more intense assertion of power until the same level of real existence is reached. 
\end{enumerate}
Then at this point:

\begin{quotex}
in correspondence with a magical act (because a \emph{magician could be defined as one who knows how to influence nature itself}), [the idea] would cease to be ``false" and ``subjective" by making itself instead ``true" and ``objective". What was error or illusion becomes accordingly truth, through power. That is valid for any empirical, moral, or rational field whatsoever. 

\end{quotex}
Evola here does not spell out the obvious corollary. Whatever happens in any ``empirical, moral or rational field whatsoever" is the result of the intention of an agent, someone or group that understands, and uses, this process or, better said, ``magicians". How much of the world is under the spell of some collective hallucination or another? Only the Trial by Fire can break the spell. Evola's goes on to demonstrate the relationship between ontology and ethics.

\begin{quotex}
Real means perfectly willed. Unreal will mean therefore privation of the Will, imperfect and insufficient Will — and this is \emph{evil}, just as, on the other hand, it is error. \emph{Evil and error are one and the same thing in the sense that the common foundation to both is impotence}. Those doctrines are in truth of a fundamental value, currently almost forgotten, that placed evil in matter, meaning by this that residue of necessity that resists the idea — that is, freedom — and limits it. Matter — the Platonic ``other" — is effectively the sign of imperfection of the act and, in as much existing, is connected to that fundamental ``injustice" about which Anaximander, Parmenides, and Empedocles speak — and it is, essentially, evil. There is no other evil beyond necessity, of which matter, brute existence, is the evidence; and as long as the ``I" will not be able avoid the experience of matter — of this ``other" against and beyond him — he will be imperfect: \emph{evil, impure, irrational}. 

\end{quotex}
It is clear that freedom is not about satisfying desires, which is really a lack, a privation. Concupiscence represents a necessity that binds and restricts our will. The purpose of the Trial by Love, by substituting the satisfaction of personal desires with a disinterested love for the genuinely good or pleasurable, is to overcome concupiscence. Thus, commands such as ``love your enemy" are to be understood in this light, as means to an end, not a universal moral principle in the Kantian sense. A blind hatred, on the other hand, is the result of an impure and irrational will, and often leads one to make poor or ineffective decisions.

Weakness of the \emph{thumos} leads a man to eschew suffering. Every form of pain or upset is felt to be an injustice. Instead, the Trial by Suffering will develop a detached attitude. In this respect, we can mention the idea of ``intentional suffering", that is, deliberately choosing the more arduous path in order to develop Power. Once again, commands such as ``turn the other cheek" must be understood in this light.

Now Evola summarizes his position on the identity of power and morals:

\begin{quotex}
So the ethical problem and the metaphysical problem coincide: the measure of reality, as of the good, of the certain and of the true, is the perfection of actuality or power. \emph{Per virtutem et potentiam idem intelligo}. [Spinoza, \emph{Ethics}, ``by virtue and power I mean the same thing."] The good or virtue is potency, evil is impotence. There is no other evil except insufficiency and weakness, no other good except the will that is unconditionally sufficient in itself, and therefore free to be and do what it wills. The rest is nonsense and can be left to women, mystics, and poets. That today the meaning of the moral as cosmic worth has been lost, that up to now, it has been reduced to nothing more than the corroboration and canonisation of deficiency, weakness, and fear — of the plague of beautiful sentiments, of noble virtues, of holy ideals — that is, of the fundamentally \emph{immoral}. 

\end{quotex}


\flrightit{Posted on 2011-03-17 by Cologero }

\begin{center}* * *\end{center}

\begin{footnotesize}\begin{sffamily}



\texttt{Albion on 2011-03-19 at 01:15 said: }

Since the purpose of the sight is not to defend Evola, but to build on him, then some elective affinity would seem to be in order, just to understand his position. However, Cologero may have stumbled on something here, which goes to the root of the misunderstanding of Platonism as somehow anti-matter. From the higher perspective, ``matter" would of course be entirely good insofar as it existed at all, but would necessarily be ``evil" insofar as it tended to deprive one of that higher perspective. Which is exactly what has happened. Plato never spoke as if we did not have bodies (George Parkin Grant). And one must not overlook the sophistication and symbolism and secret doctrine which operated to imbue statements that were ``wrong" if taken ; context, more than terminology, is everything.


\hfill

\texttt{Albion on 2011-03-19 at 12:59 said: }

I am beginning to think Platonism is nothing like what we have been told that it was. Necessity, without freedom is evil, just as freedom without necessity is evil. The traditionalists (or Platonists) don't get to talk about the divine loop where the end meets the beginning, generally because everyone quarrels with their quest. You can't have ``Western" without Platonism. I don't see where Evola is any more anti-material than Platonism, in fact, far less so.


\hfill

\texttt{Albion on 2011-03-19 at 13:00 said: }

Perhaps one could say that untransfigured matter, or matter rendered untransfigurable, is what is ``evil". Which can be done.


\hfill

\texttt{Will on 2011-03-20 at 11:25 said: }

The five numbered points from Evola's book are strongly parallel to the methodology of Tantric sadhana, so much so that I wonder whether Evola got it from Woodroffe or some other source, or whether he somehow figured it out on his own. This methodology survives to this day in Tantric Buddhist practices in which yidam deities are visualized and meditated upon and prayed to, until they become as real and more real than ordinary reality.

These visualizations are one of the primary uses of thangka paintings: they provide the image to be formed in the mind. There is some evidence that icons in the Christian tradition may have been used in similar ways, and perhaps still are in certain Orthodox circles. In both cases, the painter must go through certain preparatory and purificatory rituals prior to beginning the work. In both cases, the resultant images (if they are successful) are regarded as holy, and as retaining something of the essence of the deity or saint which they portray.


\hfill

\texttt{kadambari on 2011-03-21 at 12:22 said: }

@Cologero

``Idiot" originally referred to a ``person so mentally deficient as to be incapable of ordinary reasoning".

Alain de benoist gives an interesting definition of Idiot. I find his understanding of antiquity of interest, and do not know much about his politics. He writes:

``The notions of citizenship, liberty, or equality of political rights, as well as of popular sovereignty, were intimately interrelated. The most essential element in the notion of citizenship was someone's origin and heritage. Pericles was the ``son of Xanthippus from the deme of Cholargus." Beginning in 451 b.c., one had to be born of an Athenian mother and father in order to become a citizen. Defined by his heritage, the citizen (polítes) is opposed to idiótes, the non-citizen—a designation that quickly took on a pejorative meaning (from the notion of the rootless individual one arrived at the notion of ``idiot"). Citizenship as function derived thus from the notion of citizenship as status, which was the exclusive prerogative of birth. To be a citizen meant, in the fullest sense of the word, to have a homeland, that is, to have both a homeland and a history. One is born an Athenian—one does not become one (with rare exceptions). Furthermore, the Athenian tradition discouraged mixed marriages. Political equality, established by law, flowed from common origins that sanctioned it as well. Only birth conferred individual politeía.

That is essentially how caste was : people with a certain history and way of life…so being out of caste is to have no history and be a rootless individial…

\url{http://home.alphalink.com.au/\~{}radnat/debenoist/alain14.html}


\hfill

\texttt{John Bardis on 2011-03-21 at 16:47 said: }

In \textit{The Yoga of Power} Evola writes:

``According to Indo-Aryan tradition, a caste corresponds to an individual differentiated nature, and to the body of laws (dharma) one should follow. In the Left-Hand Path, however, dharma is seen as a limitation, or as a conditioning to be overcome…Likewise, it does not matter to what caste a person pursuing absolute liberation belongs. Obviously, this is a case of apparent `democracy'. The capability of effectively walking on the path produces differentiations among people that are more clear-cut and real than the ones due to caste membership, especially when the latter have lost their original and legitimate foundation." (page 97)

Perhaps this is what Gurdjieff's toasts to the idiots (or whatever it was called) was all about. I forget how it works, but there were forty or more different types of idiot, and each person had the task of determining which, exactly, kind of idiot he or she was–resulting in the downing of vast quatities of alcohol.


\hfill

\texttt{kadambari on 2011-03-21 at 18:57 said: }

``Likewise, it does not matter to what caste a person pursuing absolute liberation belongs."

Yes in classical India you had the sanyasin (samana) who was beyond caste and religion (you could not demand of a sanyasin what his caste traditionally). However, this did not injure the overall political set up, even the sanyasin is something at the beginning. This is what I find so amazing about this civilization: all these elements could co-exist (the samanas later lead to the rise of Buddhism), which in turn co-esists with what precedes it, so a civilization self-corrects as long as its essential character is not disturbed in the sense that the essential way of life is destroyed; all is self-correction comes to a halt when classical civilization is destroyed by outside forces in the middle ages…So this is a case of pluralism which has none of the aspects we associate with pluralism in the modern sense, it is not multicultural because the differences are a facet of a larger civilization which encompasses it all in a unity….


\hfill

\texttt{John Bardis on 2011-03-21 at 19:37 said: }

So, then, if I may ask–What caste are you?


\hfill

\texttt{kadambari on 2011-03-21 at 20:40 said: }

Well we are Brahmins and retain some memory but today the civilization there is not one that can be properly be called ``Hindu"; so I feel that we are helpless residues of a civilization that harldy exists anymore except among some remnants who carry a memory of it, people who are small in numbers; when I read this blog I feel Westerners are also dealing with similar dilemnas in a Western context…Yes I suppose we are also idiots in Benoist's sense, being unable to relate to or sympathize with the current form of civilization which has taken root there, so rootless in this sense…

I suppose there is not much one can do but our tiny part to keep the memory of these things alive…Imagine if it were not for all these traditionalist writers we would not be concerning ourselves with such questions as a larger group, even if we might do so quietly as individuals! These things have to have impact amongst the public to have a future and cannot just be confined to isolated individuals in academia…A knowledge that there is a problem is the first step…


\hfill

\texttt{John Bardis on 2011-03-22 at 12:18 said: }

In the town where I live, a few miles south of Atlanta, there's a Hindu temple. Or, actually, there are two of them side by side. They are quite magnificent. The main ``god" of the first temple has to be the most amazing statue I've ever seen.

But in another part of Atlanta there is a whole temple complex that makes our two temples seem quite small in comparison.

There was a time in my life when I had the good fortunate to attend a good many concerts of classical Indian music.

Obviously Hinduism (or whatever you call it) isn't going anywhere. It's definitely here to stay.


\hfill

\texttt{kadambari on 2011-03-22 at 14:23 said: }

I hope the Hindus are also doing something for the community… I can say that I myself know many things about our religion from the great Western scholars; the Buddhism book my brother read in college was written by a Catholic priest and it was a very good introduction I recall. Of course there are certain things that are passed down and you know as being a part of the culture, but we are always amazed by the dedicated Westerners who study it to such depths; the German scholars preserved many things when they were dying out even in our regions; genuine quest for knowledge transcends cultures…

I wish Indians would be exposed to the Greeks for they would really love them and learn so much from them and find them not different at all, but what can you do? With the current government run educational system there, Hindus barely understand even their own culture and history…


\hfill

\texttt{kadambari on 2011-03-23 at 08:31 said: }

@John Bardis

The kinds of things you talk about are probably organized by South Indians, who as classical civilization was destroyed much later in the south still have some memory of architecture and other cultrual memories and so on, and can be benign people…The North (I am speaking of places like Delhi) has a thuggish culture because it was totally destroyed at several points and hence there is a barbaric legacy which survives there, I really dislike the culture in those parts, however, the mountian/hill areas in the North where people took refuge during turmoil historically that I am familiar with are completely benign and still very pleasant to live…

As for Indian classical music, most of the great names (about ten) have died, there are only two great masters left…I wonder with interest diminishing in these things when we will again hear music composed by such masters again…


\hfill

\texttt{Yuhomo on 2011-03-23 at 09:54 said: }

Evola is more magician or philosopher than metaphysician. Here he tends to pit freedom and will against necessity, yet he forgets that freedom doesn't exist without necessity. The divine by necessity willed being and cosmos, not the fall. And to be an individual being is not a constraint or weakness since its will is to be an individual whose immanence does not run counter to transcendence. Thus the statement, ``There is no other evil beyond necessity, of which matter, brute existence, is the evidence" is absurd. Then in the moral argument Evola tends to be too Nietzschean in wanting to create an amoral morality based on power. But what do I know?


\hfill

\texttt{John Bardis on 2011-03-23 at 13:25 said: }

Yes, I was told that the temples near where I live were for south Indians.

You will have to forgive the fact that my ignorance on these matters far outweighs my very real interest.

I believe that Buddhism was still prevalent in India, though only in northern India, at the time of the Muslim invasion. And, of course, this invasion only penetrated so far. So I suppose it does make sense that the traditions would be more intact in the south.

As for the music, it seems to me quite similar to western classical music in that it could only be created and can only be performed due to the existence of child prodigies. I believe both kinds of music are still being performed at a very high level. But perhaps creation has come to a stand-still. And this seems to be the case in many fields–literature, philosophy, art. Why this is, I don't know. Perhaps there are times of creation and times of preservation. Or perhaps all the emphasis now is on technology and popular culture.


\hfill

\texttt{kadambari on 2011-03-23 at 15:46 said: }

@John Bardis

Yes the South is quite different from the North, they have their own variety of HInduism and their peculiar languages which sound even peculiar to us and their ways are quite different, but I find in spite of all these differences, the people there are gentler compared to the North; many of them obey rules as compared to the North and are capable of better organization amongst themselves, whereas the North is very unruly and disorderly owing to its bad history. Of course, there are subtleties across regions and your geography will play a part in the peculiar aesthetics of your ethnic group…

Indians unfortunately are not exposed to Western classical music that much, which is a shame because it deprives them of something amazing out there just as they are not exposed to Greco-Roman civilization so cannot properly understand Western civilization and are mostly only exposed to English in terms of language. It is not that there is no interest in these things; they just have not been exposed…There needs to be exposure to these things if Indians can have a proper perspective on civilizations not their own…

As far as Western music is concerned, most are only exposed to pop music; you see the Chinese and Japanese really into Western classical music…I was lucky to have parents who appreciated Western classical music, so was lucky to grow up appreciating both classical Indian and classical Western music…

These days, there is much exposure to books in our region which is a good thing about the internet; today you can find most of the classics on the internet which is good for areas where they are not easily accessible. My husband who is in the sciences recalls when he was in college, the librarian was a Kashmiri Pandit, so in addition to the science books the librarian had a very peculiar section of other books according to his taste such as works by Aurel Stein, the explorer Richard Burton and old British histories of India (which contained history without whitewash) etc., which he says opened him up to understanding Indian history properly, as history is not taught in the schools correctly, as it is all written with a political agenda and dominated by Marxist historians and rewriting of history to serve political ends, the truth is not presented. Most people end up studying technical subjects and the sciences to survive in the modern world and find themselves when they get older woefully ignorant of their own civilization apart from what their parents have taught them and often want a refund on their education…


\hfill

\texttt{kadambari on 2011-03-23 at 17:04 said: }

@John Bardis

We need geniuses to bring back the music, but proper patronage is also required to sustain the arts. The classical musical heritage of the South has to to with the renaissance of the arts during the Vijayanagar empire, the last Hindu stronghold. The rulers patronized the arts and all Brahmanical traditions such as music.

\url{http://www.planetradiocity.com/musicopedia/music\_decade.php?conid=2311}

\url{http://www.vasanthvisuals.com/hampi.html}

There are ragas not which have disappeared due to lack of use, many we cannot even imagine because they have been destroyed with civilization in certain parts, so what is existing cannot continue wihout proper patronage. Unfortunately, the way the Indian State is today controlled by the socialist Congress ever since after independence with the exception of five years, leading to India's current misery and ongoing deracination and overpopulation, I am skeptical of preservation of the arts unless some other party comes into power nationally as the native culture is looked down upon and not encouraged by the state…


\hfill

\texttt{Cologero on 2011-03-23 at 19:25 said: }

The ``divine" will is absolutely free and does not act out of necessity. It the divine is infinite, where would the ``necessity" come from?

Evola is speaking about trans-human states, not the ``individual" man as you seem to mean it.

Evola criticizes Nietzsche's ``morality" because the latter did not understand ``asceticism".

Please read the material carefully prior to commenting.

Evola may be wrong, but not for the reasons you present.

To emphasize once again, we are not out to convert, but rather to elucidate Evola and Guenon, in particular. The only issue at stake is whether or not we have presented Evola's philosophy clearly and accurately.


\hfill

\texttt{Yumo on 2011-03-24 at 09:18 said: }

The divine will is not absolutely free in the sense that you imply. It cannot not will the universe, which occurs as a result of its infinitude. Necessity relates to the absolute, and freedom to the infinite, both of which coincide. If the divine is infinite how could it not also include necessity? But necessity does not relate to impossibility or falsehood as Evola would imply.

I know what Evola is trying to say but his accuracy is off, and for reasons which I suspect to be his early influences, to reduce everything to the freedom of the will. So while those he criticized erred in the direction that negated transcendence, he errs in the direction that negates ethical virtue.

He says that matter is a sign of imperfection as if necessity of willing the universe was not necessary, but then the divine would not be infinite nor free. Matter is a sign of the will to matter which enabled the divine to be immanent and therefore is what it is.


\hfill

\texttt{kadambari on 2011-03-25 at 09:26 said: }

``I believe that Buddhism was still prevalent in India, though only in northern India, at the time of the Muslim invasion. And, of course, this invasion only penetrated so far."

The destruction was massive and total in most areas which never recovered. I recall the European who discovered the ellora caves. The people of the area had become so wretched, he could not imagine that the same people had once produced the art in Ellora….


\hfill

\texttt{Matt on 2011-03-25 at 16:03 said: }

Yumo, 

How does Evola negate ethics in his thought? He stated many times in his writings that he takes virtue into account, but virtue as it was traditionally defined (virtus as ``power", the power to act in the right way in the right situation). Gornahoor has a post on St. Anthony and Evola that can give you a good idea of the ethics in Evola's thought.


\hfill

\texttt{John Bardis on 2011-03-25 at 22:19 said: }

Dear Kadambari,

Thank you for those links about the Vajayanager empire. They give a context to things that I have seen and heard. These things aren't simply from India. They are from a particular place in India at a particular time.

And if these things that I have seen and heard are just the last rays of light from a star that no longer exists–well, is there anything unique in that? Here is something I happened to read today:

``All that is mortal lives in this fear of death; every new birth augments the fear by one new reason, for it augments what is mortal. Without ceasing, the womb of the indefatigable earth gives birth to what is new, each bound to die, each awaiting the day of its journey into darkness with fear and trembling."

And yet we know that this music, this art, this religious devotion, this philosophical insight, will never die.


\end{sffamily}\end{footnotesize}
